\section*{Аннотация}

\tab Производство Систем-на-Кристалле (СнК) является крайне дорогостоящим процессом. С каждым последующим этапом производства стоимость исправления ошибки только увеличивается, поэтому производители интегральных схем стремятся обнаружить их как можно раньше. Более того, производителям важно выпустить схему в продиктованные финансовыми соображениями сроки. Отсюда формируются два главных требования к верификации --- проверке корректной работы получившейся СнК: подтверждение соответствия итогового продукта заданным требованиям с высокой степенью уверенности, а также адекватные сроки самого процесса подтверждения. 

Целью данной работы является изучение актуальных методов функциональной верификации, а также проведение их сравнения по критериям применимости и качества верификации в контексте верификации кэша второго уровня GPGPU.  

В работе приводится описание формальной верификации, симуляции и метода трасс. Устанавливаются основные концепции, на которых базируются вышеупомянутые методы; на основе этого выделяются соответствующие преимущества и недостатки. Осуществляется реализация методов формальной верификации и симуляции по отношению к кэшу второго уровня GPGPU, а также сравнение этих методов по критериям полноты, скорости и применимости к поставленной задаче.


